\section{Limitations}
\label{sec:limitations}

\paragraph{Dataset size.} DRIVE contains 20 training and 20 test images.
This is standard for retinal vessel benchmarks, but it limits the statistical
power of our comparisons. Bootstrap confidence intervals (reported in
Section~\ref{sec:results_uncertainty}) partially address this, but some of
our finer-grained conclusions (e.g., the exact deferral rate at which
confidence-aware overtakes adaptive thresholding) may not generalize to
larger datasets. The 5-fold cross-validation results
(Table~\ref{tab:cv}) confirm low variance for segmentation metrics
(Dice CV $<$ 1\%) but higher variance for uncertainty-related metrics
(Unc-AUROC CV 7.4\%).

\paragraph{Single architecture.} We use a U-Net with ResNet-34 encoder
throughout. Whether the relative ordering of uncertainty methods (TTA $>$
MC Dropout) and deferral strategies (confidence-aware $>$ global) holds
for other architectures (attention U-Net, transformer-based models) is an
open question. The confidence-aware deferral strategy is architecture-agnostic
by construction, since it operates on the output uncertainty and prediction
maps. But TTA's advantage over MC Dropout could narrow with architectures
that have dropout at multiple scales or with more expressive variational
layers.

\paragraph{Binary segmentation only.} Retinal vessel segmentation is a
binary task: vessel or background. Multi-class segmentation (e.g., organ
segmentation in CT) introduces additional complexity: uncertainty over
$K > 2$ classes cannot be summarized by a single entropy value as easily,
and the decision boundary geometry is more complex. Our confidence score
$c = 2|p - 0.5|$ is specific to binary outputs. Extending to multi-class
settings would require replacing it with a margin-based or entropy-based
confidence, which is straightforward but untested here.

\paragraph{No human-in-the-loop evaluation.} We evaluate deferral quality
by measuring error reduction on held-out data with ground truth labels.
We do not evaluate whether clinicians actually make better decisions when
presented with deferral maps. The gap between computational deferral
quality and real clinical utility depends on factors we do not measure:
how review interfaces display deferred regions, how long review takes,
and whether clinicians trust and act on the deferral signal. These are
important questions for deployment, but they are outside the scope of
this computational study.

\paragraph{Geometric TTA only.} Our TTA augmentations are restricted to
the symmetry group of the square (flips and 90-degree rotations). We do
not test photometric augmentations (brightness, contrast, gamma) or elastic
deformations. Geometric augmentations have the advantage of being
interpolation-free and exactly invertible, which avoids introducing
artifacts into the uncertainty map. Photometric augmentations might capture
additional sources of prediction fragility but would also increase $K$ and
potentially introduce noise into the aggregated prediction.

\paragraph{Threshold generalization.} Deferral thresholds are fit on a
validation set from the same distribution as the test set (DRIVE). When
applied to STARE and CHASE\_DB1 (cross-dataset evaluation), the thresholds
were not re-fit. The cross-dataset results in Table~\ref{tab:cross_dataset}
are therefore conservative: re-fitting thresholds on a small labeled
subset from the target domain would likely improve deferral performance.
Whether this holds in practice depends on the availability of labeled
data at deployment time.
